%%%%%%%%%%%%%%%%%%%%%%%%%%%%%%%%%%%%%%%%%
% Motivation letter, Summer@EPFL 2027
% Deadline: Sunday 29 November 2026, 23:59 AoE. One page ideally, max two.
%%%%%%%%%%%%%%%%%%%%%%%%%%%%%%%%%%%%%%%%%
\documentclass[11pt]{article}
\usepackage[a4paper,left=0.9in,right=0.9in,top=0.8in,bottom=0.8in]{geometry}
\usepackage[T1]{fontenc}
\usepackage[utf8]{inputenc}
\usepackage{charter}
\usepackage[hidelinks]{hyperref}
\setlength{\parindent}{0pt}
\setlength{\parskip}{0.75em}
\pagestyle{empty}

\begin{document}

{\large\bfseries Osinachi Chukwujama}\\
Saint-Étienne, France \quad·\quad \href{mailto:osinachi.chukwujama@gmail.com}{osinachi.chukwujama@gmail.com} \quad·\quad \href{https://github.com/vicradon}{github.com/vicradon}

\vspace{0.5em}
Summer@EPFL Selection Committee\\
School of Computer and Communication Sciences, EPFL
\hfill \today

\vspace{0.3em}
\textbf{Application for a Summer@EPFL 2027 research internship: inverse problems in imaging}

Dear Members of the Committee,

I am applying to Summer@EPFL because I want to spend summer 2027 working on a question that has
come to organise my studies: how to recover an image from indirect, noisy measurements by combining
a physical model of how the data was formed with a learned model of what images look like. I am a
first-year student in the Erasmus Mundus Joint Master in Intelligent Photonics (iPSRS), and the
groups of Prof.\ Volkan Cevher, Prof.\ Wenzel Jakob and Prof.\ Sabine Süsstrunk sit exactly where
that question lives: optimisation for learning, differentiable physical models of light, and
computational imaging.

My path here runs through both physics and software. I graduated with First Class honours in
Mechanical Engineering from FUTO, Nigeria, then worked for several years as a software engineer and
technical lead, most recently leading the backend of an enterprise platform of nearly 1,000 files.
That work taught me to write code other people depend on: tested, reviewed and maintainable. I
returned to study because I wanted to apply that discipline to research rather than products.

On the modelling side, I have derived and implemented denoising diffusion (DDPM) and
optimal-transport conditional flow matching in PyTorch from first principles, from the ELBO and
reparameterisation through to training and sampling, and studied how they relate to score matching
and normalising flows. On the physical side, my iPSRS coursework puts me in the lab with Fourier
optics and spatial light modulators: I generated grating and Fresnel-lens phase patterns in MATLAB,
displayed them on an SLM, and recovered the gratings' spatial frequencies from the observed
diffraction. Seeing a forward model hold up on an optical bench is what made me want to invert one.
I have also taken vision models to hardware: in my final-year project I trained YOLOv5 on a dataset I
collected, exported it through ONNX and deployed it with TensorRT on a Jetson Nano, cutting
processing time by 44\%.

I would be glad to contribute to any project that brings these together, for example using
diffusion or flow-based priors for reconstruction, differentiable forward models for imaging
problems, or the optimisation methods that make such reconstructions reliable. I learn quickly in
unfamiliar codebases, and I would aim to leave the lab with a documented, reusable result rather than
a one-off experiment.

The timing also suits my programme: I study at Université Jean Monnet until January and at the
University of Eastern Finland until May 2027, and I am free for two to three months between June and
early September before my second year at Université Paris-Est Créteil, where I will specialise in
biometrics and intelligent vision. A summer at EPFL would be the bridge between the photonics I am
learning now and the computational imaging research I want to pursue.

Thank you for considering my application.

Sincerely,\\[0.4em]
Osinachi Chukwujama

\end{document}
